\section{Applying the localized oracle to the Gaussian network}
\label{sec:assembly}
\begin{proof}[Proof of Theorem~\ref{thm:main}]
At $H=1$, the localized capacity argument in Section~\ref{sec:one-hidden} proves the required raw-risk certificate with $(v,\rho)=(2,1)$; its teacher conversion below is unchanged. Assume $H\ge2$ for the program construction. The centered independent readout and bounded hidden activations give
$\E_\mu\widehat R\le2$ for every sample, width, and scale. Take $M=2$ in Section~\ref{sec:localized-direct-oracle}.

Apply Proposition~\ref{prop:programs} with the proof measure
$\xi_S=(D_X+\widehat D_{X,S})/2$. Quantize its fixed compact program family in uniform norm at accuracy $c_Ha$. After rescaling $a$ by a depth-dependent constant, this gives exactly the finite dictionaries required there, with $b=1$, $p=H^2+2H-3$, $q=4H-2$, and a dimension factor polynomial in $d+1$. The logarithmic factor in \eqref{eq:program-entropy} is absorbed by the larger explicitly chosen $\theta_n$. Thus
\[
 \rho=(q+2)/2=2H,\qquad v=\max\{p+1,2\rho\}=v_H.
\]
The same dictionaries cover every finite width and Gaussian scale. The statistical event is consequently common to all these learners; there is no union over a continuum of priors.

For a fixed teacher, Lemma~\ref{lem:teacher} and the variational formula for free energy give, for $\lambda\ge1/n$,
\[
 F_S(\lambda)\le\widehat R_S(g)+\lambda[1+2A_g(\lambda)]
 \le\widehat R_S(g)+\lambda[1+2A_g(1/n)].
\]
Also $\|g\|_\infty^2\le m^2B^2\le E_g$, so its output envelope is controlled by the comparison resource. The scalar Bernstein event for this one teacher is common to all learner widths and scales. Allocate half the failure probability to it and half to the fixed-dictionary event. The population-oracle argument of \eqref{eq:localized-population-oracle}, with the enlarged resource $\mathcal A_{g,n}$, proves \eqref{eq:main-oracle}. The fallback branch is covered as in that argument.

Every finite posterior has finite readout moments because its density is bounded by $Z^{-1}$ relative to its Gaussian prior. Jensen's inequality therefore proves the same bound for the ordinary raw posterior mean.
\end{proof}

\begin{proof}[Proof of Corollary~\ref{cor:rates}]
At fixed resources, $\mathcal A_{g,n}=O_H(\log(n/\delta))$ and $\theta_n$ is $n^{-1}$ times logarithmic factors. This proves the two rates. For $m\le n^\alpha$ and $B\le n^\beta$, at fixed $d,H$, the resource is at most $O_H(n^{2(\alpha+\beta)}\log n)$. The first remainder vanishes when $2(\alpha+\beta)<1/v_H$; the second gives the same inequality because $r_H=v_H/\rho_H$.

A minimum-risk teacher exists in each compact bounded-weight class by dominated convergence. The zero-readout network belongs to that class, so its optimal risk is at most one. Apply the theorem to this deterministic population minimizer. The calibration does not depend on the teacher class. A countable family of comparison classes may be covered by summable confidence allocations, adding only the corresponding logarithmic confidence costs.
\end{proof}

\paragraph{Measurability.}
Finite program families are compact continuous images of bounded parameter boxes and have countable uniformly dense subsets. Actual-network norm tests and residual tests are Borel in the finite weights and sample; spectral bases can be chosen measurably from finite Gram matrices. Population integrals are measurable by bounded convergence. The statistical suprema are over fixed dictionaries, whereas sample-dependent good sets enter only prior integrals whose bounds hold for every fixed sample. No proof-program selection is part of the learner.
